\originalchapter{恒星种群、形成史与化学守恒}

本章对应 8.902 第8讲和第9讲开头的 G 矮星问题. 前面把恒星视作有质量的轨道示踪体, 现在恢复它们的出生、寿命、光谱和元素回馈. 光谱合成把形成史映入观测, 化学演化则用质量守恒约束形成史; 两种信息可以互补, 却都需要初始质量函数、混合和流入流出的假设.

\originalsection{形成率的指示量不是无条件换算表}
恒星形成率 $\psi(t)$ 是单位时间新形成恒星的总初始质量, 常用 $\Ms\,\mathrm{yr^{-1}}$. 年轻大质量恒星产生较强紫外与电离光子, 后者使 H II 区产生 H$\alpha$; 尘埃吸收部分光再辐射于远红外. 因而这些光度能作为形成率指示量, 但响应的时间窗、尘埃和 IMF 会影响标定.

课程采用的数量级关系为
\begin{align*}
\frac{\psi_{\rm FIR}}{\Ms\,\mathrm{yr^{-1}}}&\sim\frac{L_{\rm FIR}}{5.8\times10^9\Ls},\\
\frac{\psi_{\rm H\alpha}}{\Ms\,\mathrm{yr^{-1}}}&\sim\frac{L_{\rm H\alpha}}{1.3\times10^{41}\,\mathrm{erg\,s^{-1}}},\\
\frac{\psi_{\rm UV}}{\Ms\,\mathrm{yr^{-1}}}&\sim\frac{L_{\nu,\rm UV}}{7.2\times10^{27}\,\mathrm{erg\,s^{-1}Hz^{-1}}}.
\end{align*}
UV 项须使用相应的频率光度密度; 原笔记只写 $\mathrm{erg/s}$ 会使量纲与该标定不符. 这些数是课程在给定模型背景下的近似例子, 不是所有金属丰度、突发形成史和 IMF 通用的常数. 老恒星或 AGN 也能加热尘埃, 吸收又可隐藏紫外和 H$\alpha$, 应比较不同指示量而非任取一个光度直接宣称唯一形成率.

冷而密的气体可在自由落体时间内聚集, 常用局部尺度模型 $\psi\sim\epsilon_{\rm ff}M_g/t_{\rm ff}$, $t_{\rm ff}\propto(G\rho)^{-1/2}$. 系数 $\epsilon_{\rm ff}$ 包含湍流、磁场和反馈等未由这一式求出的物理. 分子相比例还依赖屏蔽、金属丰度和团块结构. 原课受限论文的气体相图和观测散点图不重新分发; 这里解释模型依赖, 不把相图百分数当作普适宇宙常数.

\originalsection{IMF 归一化决定每个积分的单位}
用物理质量 $m$ 定义 $\phi(m)$, 使总初始质量 $M_{\rm new}$ 的一次形成产生
\[
\dd N=M_{\rm new}\phi(m)\dd m,\qquad
\int_{m_l}^{m_h}m\phi(m)\dd m=1.
\]
于是 $[\phi]=M^{-2}$, 右侧归一化是无量纲1. 若改以 $x=m/\Ms$ 表示, 可定义无量纲 $\varphi(x)$ 使 $\int x\varphi(x)\dd x=1$, 则 $\dd N=(M_{\rm new}/\Ms)\varphi(x)\dd x$. 原课把这两种单位约定交叉书写, 本卷统一使用前一种.

Salpeter 形式 $\phi=Cm^{-2.35}$ 在选定 $m_l,m_h$ 上由上述质量积分定 $C$. 同一幂律下数目与质量贡献分别按 $m^{-2.35}\dd m$ 和 $m^{-1.35}\dd m$ 加权. 低质量恒星多且长寿, 高质量恒星少却可能主导紫外光, 因而按光度选样会偏向后者. 课程的 $m_l\sim0.08\Ms$ 对应氢燃烧门槛尺度, $m_h\sim100\Ms$ 是模型截断; 不能只由 $L_{\rm Edd}\propto M$ 宣称100太阳质量是所有恒星绝不能超过的硬上限.

\originalsection{寿命与合成光谱的双重积分}
恒星的有效温度按 $L=4\pi R^2\sigma_{\rm SB}T_{\rm eff}^4$ 定义, 不要求真实频谱在每个频率上都是理想黑体. HR 图通常横轴温度向右降低, 纵轴光度增加. 固定半径线满足 $\log L=4\log T+\mathrm{const}$, 但主序的物理定义是核心氢燃烧, 各主序恒星半径不同. 原课将 $R=1$ 线直接叫主序, 此处纠正这一概念.

假定可燃料量与质量近似成正比, 则主序寿命 $\tau_{\rm MS}\propto M/L$. 在采用 $L\propto M^p$ 的质量区间, 有 $\tau_{\rm MS}\propto M^{1-p}$. 课程以 $p\sim4$ 得 $M^{-3}$, 低质量例子取 $p\sim3$ 得 $M^{-2}$. 这些是区间近似, 不可外推整个质量谱. 主序转折质量由 $\tau_{
\rm MS}(m_t)\simeq$ 种群年龄给出, 为老种群年龄提供线索.

设单星频率光度为 $\ell_\nu(m,Z,\tau)$, 对同年龄同初始金属丰度的种群, 每单位初始形成质量的光谱为
\[
\ell_\nu^{\rm SSP}(\tau,Z)=\int\phi(m)\ell_\nu(m,Z,\tau)\dd m.
\]
单星演化和残骸阶段含在 $\ell_\nu$ 内, 再叠加出生时刻,
\[
L_\nu(t)=\int_0^t\psi(t')\ell_\nu^{\rm SSP}(t-t',Z(t'))\dd t'.
\]
先选 IMF 和形成史, 再生成光谱、通过尘埃与滤镜映入数据, 比较后调整, 才是合成的工作流. 年龄、金属丰度、尘埃和 IMF 的效应可退化, 线特征与动力学质光比能提供独立限制. 本卷没有假称已运行某个完整种群合成软件或拟合真实星系全谱.

\originalsection{元素丰度与四条质量守恒式}
金属质量分数 $Z=M_{\rm heavy}/M$ 与元素数目比的对数
\[
[X/Y]=\log_{10}\frac{(n_X/n_Y)}{(n_X/n_Y)_\odot}
\]
不同. $[\mathrm{Fe/H}]=-1$ 表示铁氢数目比为相应太阳标定的十分之一, 不直接意味着所有重元素总质量分数都等于太阳的十分之一. 课程的 $Z_\odot\simeq0.02$ 是所用近似标度; 不同丰度标定需显式说明.

令气体质量 $M_g$, 恒星及残骸中锁定质量 $M_*$, 总重子质量 $M=M_g+M_*$. 形成率为 $\psi$, 恒星返还气体率为 $E$, 返还金属率为 $E_Z$, 流入率为 $F$、金属分数 $Z_F$, 流出率为 $O$、金属分数 $Z_O$. 均匀混合模型有
\begin{align*}
\dot M&=F-O, &\dot M_*&=\psi-E,\\
\dot M_g&=-\psi+E+F-O,
&\frac{\dd(ZM_g)}{\dd t}&=-Z\psi+E_Z+Z_FF-Z_OO.
\end{align*}
其中恒星形成先锁定原有金属, 返还包含原有金属和新核合成金属. 不能只保留新金属而漏掉原有返还部分. $Z_O$ 可不同于平均 $Z$, 选择性金属风与完全混合风的结果不同.

若质量 $m$ 的星在 $\tau_m$ 后返还 $m-w_m$, 则
\[
E(t)=\int(m-w_m)\psi(t-\tau_m)\phi(m)\dd m.
\]
积分只计当时已演化的恒星, 出生前取 $\psi=0$. 相应金属率可写
\[
E_Z(t)=\int\![(m-w_m)Z(t-\tau_m)+p_Z(m)]
\psi(t-\tau_m)\phi(m)\dd m,
\]
其中 $p_Z(m)$ 是新产生并释放的金属质量, 单个元素的净产生量也可能为负. 延迟时间使金属史并非仅由当前 $\psi(t)$ 决定.

\originalsection{瞬时回收近似与混合流出的抵消}
对寿命远短于气体演化时间的返还星, 近似 $\psi(t-\tau_m)\simeq\psi(t)$ 并立即混合. 定义返还比例和按锁定质量计的净产额
\[
\mathcal R=\int_{m_{\rm lim}}^{m_h}(m-w_m)\phi(m)\dd m,
\qquad y=\frac1{1-\mathcal R}\int_{m_{\rm lim}}^{m_h}p_Z(m)\phi(m)\dd m.
\]
于是 $E=\mathcal R\psi$, $E_Z=\mathcal R Z\psi+(1-\mathcal R)y\psi$. 注意锁定质量包括长寿恒星和残骸, 不只是残骸. 用 $\dd(ZM_g)/\dd t=M_g\dot Z+Z\dot M_g$, 两条守恒式相减得到
\[
M_g\dot Z=(1-\mathcal R)y\psi+(Z_F-Z)F-(Z_O-Z)O.
\]
若流出完全混合, $Z_O=Z$, 流出项在金属分数的瞬时导数中抵消; 它仍通过减少 $M_g$ 改变演化. 原讲义相应公式保留 $+Ze$ 项, 与前面的混合流出守恒式不一致. 不能因为符号展开失误把流出误认为直接提高金属分数.

瞬时回收对低质量长寿星和 Ia 型超新星延迟等不适用, 不应使用该近似解释所有元素比的时钟. 后面解析闭盒结果先限制在它的适用模型中.

\originalsection{闭盒解与恒星的平均金属丰度}
取 $F=O=0$, 初始全为无金属气体 $M_g(0)=M$, 则 $M$ 恒定,
\[
\dot M_g=-(1-\mathcal R)\psi,
\qquad M_g\dot Z=(1-\mathcal R)y\psi.
\]
消去时间得到 $\dd Z=-y\dd M_g/M_g$, 故以气体分数 $u=M_g/M$ 表示,
\[
Z=y\ln(1/u),\qquad M_*=(1-u)M.
\]
形成率决定到达某个 $u$ 需要多久, 却不改变这一 $Z(u)$ 的闭盒关系. 所有产生的金属在气体和锁定恒星中, 总金属质量为 $yM_*$. 因而恒星质量加权平均
\[
\langle Z_*\rangle M_*+ZM_g=yM_*,\qquad
\langle Z_*\rangle=y-Z\frac{u}{1-u}.
\]
当 $u\to0$ 平均趋 $y$, 而气体 $Z$ 形式上发散. 这只是极少量剩余气体中积累金属的模型极限, 到金属分数不再小或产额假设失效时不能继续物理外推. 当 $u\to1$ 时两项相消给初始平均接近零, 不是负金属丰度.

\originalsection{G 矮星问题约束的不只是一个平均值}
长寿恒星保留出生时的气体丰度. 闭盒单调增金属的模型中, 丰度不超过 $Z$ 时已锁定的质量为 $M[1-\exp(-Z/y)]$. 若现在气体分数为 $u_0$, 现有气体丰度 $Z_0=y\ln(1/u_0)$, 则
\[
\frac{M_*(\le Z)}{M_{*,0}}=
\frac{1-\exp(-Z/y)}{1-u_0},\qquad 0\le Z\le Z_0.
\]
微分分布正比 $\exp(-Z/y)$, 一直预测相当份额的低金属长寿星. 在恒定 IMF、样本寿命选择等条件下, 才能把这个质量比例映成某类 G 矮星的数目比例. 课程以实际低金属星较少的情况指出闭盒模型缺口, 候选修正包括流入、预增丰度与形成史假设; 不能仅靠平均 $\langle Z_*\rangle$ 替代整个分布检验.

作为可检验的模型扩展, 若完全混合流出 $O=\lambda\psi$, 无流入, 常 $\lambda$, 得
\[
Z=y_{\rm eff}\ln\frac{M_{g,0}}{M_g},\qquad
y_{\rm eff}=y\frac{1-\mathcal R}{1-\mathcal R+\lambda}.
\]
这说明混合流出虽在 $\dot Z$ 显式项中抵消, 却降低按剩余气体比推断的有效产额. 若有无金属流入, 则 $-ZF$ 直接稀释, 又是另一机制. 这些有限模型从守恒式推导, 不是假称已唯一拟合银河化学演化. 下一章的 AGN 将增加能量与反馈信息, 与本章质量、金属预算相联系.
